\chapter{Experimental Setup}\label{app:experimental-setup}

This appendix records the configuration that produced every empirical
result reported in Chapter~\ref{ch:experiments}. The content
complements the high-level methodology of
Section~\ref{sec:experimental-design} with the per-run hyperparameter
schedule, the ASR baseline architecture, the dataset-construction
details, the hardware and software stack, and the reproducibility
artefacts associated with each cell of the experimental matrix.

\section{Training configuration}\label{app:setup:training}

The training schedule is the locked 50-epoch matched budget shared by
every LibriSpeech and Common Voice cell of the matrix. The schedule
was set after an early discovery phase on a 30-epoch budget and frozen
once the matrix proper began. Per-cell \texttt{config.yaml} files
under \texttt{experiments/final/outputs/} carry the snapshot that
produced each run; the values in
Table~\ref{tab:setup:hyperparams} are reproduced verbatim from those
snapshots and were verified against three independent cells (one per
architecture).

\begin{table}[h]
\centering
\small
\caption[Locked training hyperparameters]{Locked training
hyperparameters used by every LibriSpeech and Common Voice cell in the
matrix. The values are the same across architectures and scales unless
explicitly noted.}
\label{tab:setup:hyperparams}
\setlength{\tabcolsep}{6pt}
\renewcommand{\arraystretch}{1.15}
\begin{tabularx}{\linewidth}{@{}p{4.5cm} X@{}}
\toprule
\textbf{Hyperparameter} & \textbf{Value} \\
\midrule
Optimiser & AdamW \\
Peak learning rate & $3 \times 10^{-4}$ \\
Weight decay & $0.01$ \\
Learning-rate schedule & Cosine decay with linear warmup \\
Warmup steps & $1000$ \\
Epochs (matched budget) & $50$ \\
Early-stopping patience & $15$ epochs on dev CER \\
Batch budget & $\le 300$ s of audio per batch \\
Gradient clipping ($\ell_{2}$ norm) & $5.0$ \\
SpecAugment policy & LD: freq mask $F = 27$, time mask $T = 100$ \\
SpecAugment masks per utterance & 2 frequency masks, 2 time masks \\
Dropout & $0.1$ \\
Random seed & $42$ (single-seed default) \\
Numerical precision & fp32 baseline; bf16 only for the LION-mode
LA cells where memory pressure required it \\
Mamba-2 SSD chunk size & $64$ \\
\bottomrule
\end{tabularx}
\end{table}

The cosine schedule is applied to the learning rate after the
$1000$-step linear warmup completes; the schedule decays to a small
positive floor over the remaining steps of the 50-epoch budget. The
patience-15 early-stopping rule is set against dev CER and was not
triggered for any reported cell, since the matched 50-epoch budget
sits within the convergence horizon of all three backbones at both
scales.

The batch budget is duration-based rather than utterance-based. Each
training mini-batch packs utterances until the cumulative duration
reaches $300$~s, which keeps GPU occupancy stable across the
short-to-medium audio length distribution of LibriSpeech and Common
Voice without requiring per-step length statistics. Validation
batches use a fixed duration limit of $300$~s as well, with no
SpecAugment.

The single-seed convention is the matrix-wide default. Multi-seed
evaluation was reserved in advance for cells whose claim sits in the
BREAK band ($|\Delta\,\text{CER}| > 5\sigma$ vs the matched vanilla,
where the noise floor on the codebase is $\sigma \approx 0.0014$ as
estimated during the discovery phase). The matrix-as-artefact framing
of Chapter~\ref{ch:experiments} relies on the single coherent
direction of $\approx 70$ statistical tests rather than on any
single-cell multi-seed claim.

\section{ASR baseline architecture}\label{app:setup:arch}

The ASR model is identical across cells except for the backbone
dispatch. The encoder factory of
\texttt{experiments/formal\_v1/src/models/encoder.py} instantiates a
shared frontend and CTC head on top of the per-cell time-mixing
backbone; the structural pipeline is summarised in
Table~\ref{tab:setup:arch}.

\begin{table}[h]
\centering
\small
\caption[ASR baseline architecture]{ASR baseline pipeline shared by
every reported cell. The backbone row is the only stage that varies
across cells.}
\label{tab:setup:arch}
\setlength{\tabcolsep}{6pt}
\renewcommand{\arraystretch}{1.15}
\begin{tabularx}{\linewidth}{@{}p{3.5cm} X@{}}
\toprule
\textbf{Stage} & \textbf{Configuration} \\
\midrule
Audio frontend & Log-mel spectrogram, $80$ mel bins, $25$~ms window,
$10$~ms hop, sample rate $16$~kHz \\
\addlinespace
Subsampling & $2\times$ acoustic subsampling before the encoder
backbone, applied via a small convolutional stack \\
\addlinespace
Encoder block (per layer) & Time-mix block (backbone-specific) plus a
channel-mix block, each with residual connection and pre-LayerNorm \\
\addlinespace
Backbone & Per-cell choice of causal RWKV-6, Mamba-2, or Linear
Attention; or the LION bidirectional adaptation of the same three \\
\addlinespace
CTC head & Linear projection to vocabulary followed by softmax and
the connectionist temporal classification
loss~\cite{graves2006ctc} \\
\addlinespace
Decoding & Greedy CTC, no language model rescoring, no beam search \\
\bottomrule
\end{tabularx}
\end{table}

The encoder configuration at the 7M scale is $d_{\text{model}} = 256$,
$L = 6$ layers, $H = 4$ heads, head dimension $64$. At the 14M
scale the layer count doubles to $L = 12$ with the other widths
unchanged. The Mamba-2 backbone uses state dimension
$d_{\text{state}} = 64$ and SSD chunk size $T_{c} = 64$ throughout.
Total parameter counts per cell range from $6.26$~M for vanilla LA at
7M to $13.79$~M for the RWKV-6 DHO cell at 14M; per-cell counts are
reported in Table~\ref{tab:results:master-matrix-causal-7m} of
Appendix~\ref{app:results}.

The LION wrapper preserves the time-mix and channel-mix
abstraction~\cite{afzal2025linear}, replacing the causal scan kernel
with a parallel $T \times T$ attention computation under a symmetric
decay mask. The decay-class mapping per backbone is described in
Section~\ref{sec:lion} and Section~\ref{sec:proposed:lion}; the
controlled LION-S deployment on Linear Attention is implemented as a
distinct encoder class, \texttt{LIONLinearAttentionEncoder} with
\texttt{decay\_mode \(\in\) \{lit, s\}}.

\section{Datasets}\label{app:setup:datasets}

Three datasets carry the empirical evaluation, chosen for the
expressivity axis each one exercises rather than for aggregate-score
competitiveness.

\subsection{LibriSpeech clean-100}\label{app:setup:librispeech}

The principal probe is the train-clean-100 subset of
LibriSpeech~\cite{panayotov2015librispeech}: $100$ hours of
audiobook-source read English speech recorded at $16$~kHz, with the
official train, dev-clean, and test-clean splits used unchanged. The
character-level vocabulary contains the $26$ lowercase letters, the
apostrophe, the space character, and a dedicated CTC blank token, for
a head dimension of $29$.

Audio is preprocessed with the parameters of
Table~\ref{tab:setup:arch}: $80$-mel log-spectrograms with $25$~ms
window and $10$~ms hop, followed by the $2\times$ acoustic
subsampling described above. Utterances longer than $20$~s and
shorter than $0.5$~s are filtered from the training stream;
SpecAugment LD is applied during training as specified in
Section~\ref{app:setup:training} and disabled during evaluation. Test
evaluation uses full-utterance greedy CTC decoding with no language
model rescoring.

\subsection{Common Voice English 100h}\label{app:setup:cv}

The cross-distribution probe is a $100$-hour balanced subset of the
Common Voice English corpus~\cite{ardila2019commonvoice}. The
training subset is constructed from the validated split of the Common
Voice release: utterances are sampled stratified by speaker so that
no speaker appears in more than one of the train, dev, or test
partitions, and so that the cumulative duration matches the
LibriSpeech train-clean-100 budget within rounding.

The training manifest used for every Common Voice pilot cell is
\texttt{100h\_train\_manifest\_seed42.csv}, which contains $64\,198$
clips totalling $100.0005$ hours of audio. The manifest is identified
by the SHA-256 hash prefix \texttt{23f70d1c}. The dev and test
partitions follow the released Common Voice splits restricted to the
sampled speaker set. Audio preprocessing matches the LibriSpeech
pipeline; the only change is the vocabulary, which retains the same
character set since both corpora are English.

\subsection{Multi-Query Associative Recall}\label{app:setup:mqar}

The synthetic axis-2 probe is the multi-query associative recall
benchmark of~\citet{arora2023zoology}, generated from the reference
implementation. Each example is a sequence of $(k, v)$ pairs followed
by query keys whose values must be retrieved. Sequence length $T$
sweeps the canonical $\{64, 256, 1024\}$ values; the number of stored
pairs $K$ scales with $T$ such that the difficulty parameter $K/d$
stays meaningful at every length. For $T = 64$, $K = 16$; for
$T = 256$, $K = 64$; for $T = 1024$, $K = 256$. The vocabulary
contains $8192$ key tokens and $8192$ value tokens drawn from disjoint
ranges. Training uses unlimited synthetic data with cross-entropy
loss on the value-prediction positions; the success criterion is the
per-sequence accuracy reaching $0.9$. Each MQAR cell is run for up to
$30\,000$ steps with early stopping when per-sequence accuracy
reaches $0.999$ on a held-out validation batch.

\section{Hardware and software environment}\label{app:setup:hw}

All cells were trained on a single workstation equipped with NVIDIA
RTX PRO 6000 Blackwell GPUs (48~GB of VRAM each). Single-GPU
training is the matrix-wide default. The only exception is the LA
LION-S DHO cell at 7M, which used 4-GPU data-parallel distributed
training to fit the bidirectional complex-attention scan within a
practical wall-clock window; the gradient accumulation pattern was
set to keep the effective batch budget at the matrix-wide
$300$~s of audio.

The software stack consists of PyTorch with CUDA acceleration. The
RWKV-6 time-mix and channel-mix blocks are adapted from the official
RWKV-block reference implementation~\cite{rwkv2024block}; the
Mamba-2 selective state-space duality kernel is sourced from the
\texttt{mamba-ssm} package distributed by Gu and
Dao~\cite{gu2024mambacode}; the linear-attention causal recurrence is
adapted from the reference associated with~\citet{katharopoulos2020transformers}.
Modifications introduced by the proposed mechanisms are confined to
dedicated extension points described in Appendix~\ref{app:mechanisms}
and do not alter the base scan implementations, preserving numerical
equivalence with the reference at the zero-mechanism limit.

The repository commit hash that produced each run is recorded in the
per-cell \texttt{git\_sha.txt} file under
\texttt{experiments/final/outputs/<cell>/}; the training command-line
arguments are stored in the corresponding \texttt{cli\_args.txt}.
Together, these two files allow a re-run to be initiated from a
clean clone with bit-for-bit identical configuration, modulo any
non-determinism introduced by the underlying CUDA kernels.

\section{Reproducibility}\label{app:setup:repro}

Every reported cell writes a self-contained directory under
\texttt{experiments/final/outputs/<cell>\_seed42/}. The mandatory
artefacts written by the training pipeline are listed in
Table~\ref{tab:setup:per-cell-artefacts}.

\begin{table}[h]
\centering
\small
\caption[Per-cell output artefacts]{Per-cell output artefacts. Every
file is mandatory for a cell to count as reported, with the exception
of \texttt{best\_model.pt}, whose inclusion is governed by repository
storage policy and which is committed for the final-stage cells under
\texttt{experiments/final/}.}
\label{tab:setup:per-cell-artefacts}
\setlength{\tabcolsep}{6pt}
\renewcommand{\arraystretch}{1.10}
\begin{tabularx}{\linewidth}{@{}p{3.6cm} X@{}}
\toprule
\textbf{Filename} & \textbf{Content} \\
\midrule
\texttt{config.yaml} &
Full \texttt{ExperimentConfig} snapshot of the run, suitable for
direct replay. \\
\addlinespace
\texttt{cli\_args.txt} &
Command-line arguments that produced the run, including the backbone
identifier, scale, and seed. \\
\addlinespace
\texttt{git\_sha.txt} &
Repository commit hash at training time. \\
\addlinespace
\texttt{history.csv} &
One row per epoch with train and dev loss, dev CER and WER,
per-utterance CER percentiles ($p50$, $p90$, $p99$), gradient norms,
parameter $\ell_{2}$ norm, and wall-clock time. \\
\addlinespace
\texttt{metrics.jsonl} &
Streaming per-step record sampled at every $100$ training steps,
with loss, learning rate, gradient norm, and peak GPU memory. \\
\addlinespace
\texttt{results.json} &
Final test-set evaluation in full-utterance mode plus the
chunked-streaming evaluation block at $\{2.0, 5.0, 10.0\}$~s window
sizes under both reset and carry policies. \\
\addlinespace
\texttt{best\_model.pt} &
PyTorch checkpoint at best dev CER. Final-stage cells commit this
artefact under \texttt{experiments/final/} for thesis-grade
reproducibility; per-epoch diagnostic snapshots remain local. \\
\addlinespace
\texttt{plots/} &
Per-cell training-curve renderings produced at the close of each
run. \\
\addlinespace
\texttt{train.log} &
Captured stdout and stderr from the training process. \\
\bottomrule
\end{tabularx}
\end{table}

The repository is organised into four working areas:
\texttt{experiments/final/} contains the canonical matrix and the
locked plan; \texttt{experiments/formal\_v1/} contains the ASR
codebase including the encoder factory and per-backbone time-mixing
implementations; \texttt{experiments/synthetics\_v1/} contains the
MQAR codebase; and \texttt{papers/} archives the external references
consulted during the project. The reference implementations adapted
into the codebase are catalogued in
Section~\ref{sec:related:references} of the related-work chapter.

\section{Notes on cells excluded from the reported matrix}\label{app:setup:exclusions}

Two classes of cell are present in
\texttt{experiments/final/outputs/} but are excluded from the reported
matrix tables in Appendix~\ref{app:results}. The first class consists
of pre-fix and smoke-test artefacts whose directory names carry the
suffixes \texttt{\_5ep} or \texttt{\_PRE\_*}; these were retained on
disk for engineering audit but do not satisfy the locked
$50$-epoch budget specification of
Section~\ref{app:setup:training}. The second class consists of
operator-level bidirectional alternatives (\texttt{*\_bidir\_vim\_*}
and \texttt{*\_biwkv\_*}) that fall outside the unified LION wrapper
scope of Section~\ref{sec:proposed:lion}; these were retained for
internal comparison but are not part of the LION-mode rows of the
reported matrix.

The 30-epoch discovery-phase results, including the original
\texttt{rwkv6\_rse\_strong\_viscosity} anchor and the cross-stage
analyses summarised in
Section~\ref{sec:closed-cells-engaged-null}, live under
\texttt{experiments/formal\_v1/outputs/}; they are referenced by
Appendix~\ref{app:engaged-null} for the closed-cell catalogue and
are not part of the matched-budget matrix.
